\documentclass[10pt,a4paper]{amsart}
\usepackage{amsmath,amssymb,amsthm,graphicx}
\usepackage{amsmath,amssymb,graphicx,xcolor,mathtools}
\title{Unrestricted unification type of EL is infinitary, contrasted with restricted type zero}
\author{Franz Baader, Oliver Fernández Gil}
\date{doi:10.4230/LIPIcs.FSCD.2025.8; CC BY4.0}
\begin{document}
\maketitle
\noindent\textit{Editorial note.} Source https://drops.dagstuhl.de/entities/document/10.4230/LIPIcs.FSCD.2025.8. © Franz Baader and Oliver Fernández Gil, CC BY4.0. Editable text transcribed from official HTML; original publisher PDF retained.\par\medskip
\noindent\textit{Allowed standard prerequisites (editorial).}\begin{itemize}
\item Equational unification preorders, substitution permutations and reduced EL descriptions. Full author §2 definitions/proofs; §3 Theorem9 and Corollary10 cite standard EL normal-form characterization from author reference13.
\end{itemize}\medskip
% BEGIN S2
 
\par\medskip\noindent\textbf{2 Basic definitions and general results}\par
 
% BEGIN S2.p1
 
Given a signature $\Sigma$ consisting of a finite set of function symbols (with associated arities) and a countably infinite set of variables $V$, the set $T(\Sigma,V)$ of terms over $\Sigma$ with variables in $V$ is defined in the usual way [16]. An equational theory $E$ is given by a finite set of identities $s\approx t$ between terms, which are (implicitly) assumed to be universally quantified. Such a set of identities $E$ induces the congruence relation $\approx_{E}$ on terms, which can either be defined syntactically through rewriting or semantically through first-order interpretations of $\Sigma$, with $\approx$ as identity relation [16].\par
 
% END S2.p1
 
% BEGIN S2.p2
 
A substitution $\sigma$ is a mapping from $V$ to $T(\Sigma,V)$ that has a finite domain $\operatorname{Dom}(\sigma):=\{x\in V\mid\sigma(x)\neq x\}$. It can be homomorphically extended to a mapping from $T(\Sigma,V)$ to $T(\Sigma,V)$ by defining $\sigma(f(t_{1},\ldots,t_{n})):=f(\sigma(t_{1}),\ldots,\sigma(t_{n}))$. The variable range $\operatorname{VRan}(\sigma)$ of $\sigma$ consists of the set of variables occurring in the terms $\sigma(x)$ for $x\in\operatorname{Dom}(\sigma)$. Substitutions can be compared using the instantiation preorder: given an equational theory $E$, a set of variables $X\subseteq V$, and two substitutions $\sigma,\tau$, we say that $\sigma$ is more general than $\tau$ (or $\tau$ is an instance of $\sigma$) w.r.t. $E$ and $X$ (written $\sigma\leq_{E}^{X}\tau$) if there is a substitution $\lambda$ such that $\lambda\sigma\approx_{E}^{X}\tau$, i.e., $\lambda(\sigma(x))\approx_{E}\tau(x)$ holds for all $x\in X$. In case $X=V$ we also write $\lambda\sigma\approx_{E}\tau$ in place of $\lambda\sigma\approx_{E}^{V}\tau$. It is easy to see that $\leq_{E}^{X}$ is indeed a preorder, i.e., reflexive and transitive, but in general not antisymmetric. We write $\sim_{E}^{X}$ for the equivalence relation induced by $\leq_{E}^{X}$, i.e., $\sigma\sim_{E}^{X}\tau$ iff $\sigma\leq^{X}_{E}\tau$ and $\tau\leq^{X}_{E}\sigma$. We say that $\sigma$ is strictly more general than $\tau$ (or $\tau$ is a strict instance of $\sigma$) w.r.t. $E$ and $X$ (written $\sigma<_{E}^{X}\tau$) if $\sigma\leq_{E}^{X}\tau$ and $\sigma\not\sim_{E}^{X}\tau$.\par
 
% END S2.p2
 
% BEGIN S2.p3
 
An $E$-unification problem is a finite set of equations of the form $\Gamma=\{s_{1}\approx_{E}^{?}t_{1},\ldots,s_{n}\approx_{E}^{?}t_{n}\}$ such that $s_{1},t_{1},\ldots,s_{n},t_{n}$ are terms in $T(\Sigma,V)$. An $E$-unifier of $\Gamma$ is a substitution $\sigma$ that solves all the equations in $\Gamma$, i.e., satisfies $\sigma(s_{i})\approx_{E}\sigma(t_{i})$ for all $i,1\leq i\leq n$. The unification problem $\Gamma$ is solvable if it has an $E$-unifier. The set of all $E$-unifiers of $\Gamma$ is denoted as $\mathcal{U}_{E}(\Gamma)$. For elementary $E$-unification it is assumed that $\Sigma$ (and thus also $\Gamma$) contains only function symbols occurring in $E$. For $E$-unification with constants $\Sigma$ and $\Gamma$ may contain additional constant symbols, and for general $E$-unification $\Sigma$ and $\Gamma$ may contain additional function symbols of arbitrary arity.\par
 
% END S2.p3
 
% BEGIN S2.p4
 
Unification types for non-empty sets of identities $E$ are usually defined w.r.t. the restricted instantiation preorder, which is $\leq_{E}^{X}$ where $X$ is the finite set $\operatorname{Var}(\Gamma)$ of all variables occurring in the given unification problem $\Gamma$, but some authors also use the unrestricted instantiation preorder $\leq_{E}^{V}$. In this section, we will additionally consider settings where $X$ is between these two extremes. Note that $\operatorname{Var}(\Gamma)\subseteq X$ is required for the set of $E$-unifiers to be closed under instantiation.\par
 
% END S2.p4
 
% BEGIN Thmtheorem1
 
\par\medskip\noindent\textbf{Lemma 1.}\par
 
% BEGIN Thmtheorem1.p1
 
If $\Gamma$ is an $E$-unification problem and $X\subseteq V$ a set of variables satisfying $\operatorname{Var}(\Gamma)\subseteq X$, then $\sigma\in\mathcal{U}_{E}(\Gamma)$ implies $\theta\in\mathcal{U}_{E}(\Gamma)$ for all substitutions $\theta$ such that $\sigma\leq_{E}^{X}\theta$.\par
 
% END Thmtheorem1.p1
 
% END Thmtheorem1
 
% BEGIN S2.p5
 
Given an $E$-unification problem $\Gamma$ and some set of variables $X$ with $\operatorname{Var}(\Gamma)\subseteq X\subseteq V$, we say that a set $\mathcal{S}$ of substitutions is a complete set of $E$-unifiers of $\Gamma$ w.r.t. $\leq_{E}^{X}$ if it consists of $E$-unifiers of $\Gamma$, and every $E$-unifier of $\Gamma$ is an instance of an element of the complete set, i.e., for every $\theta\in\mathcal{U}_{E}(\Gamma)$ there exists $\sigma\in\mathcal{S}$ such that $\sigma\leq_{E}^{X}\theta$. Such a set is called minimal if it does not contain two distinct elements that are comparable w.r.t. $\leq_{E}^{X}$. It is easy to see that minimal complete sets of $E$-unifiers of a given unification problem $\Gamma$ are unique up to the equivalence relation $\sim_{E}^{X}$ induced by the preorder $\leq_{E}^{X}$ (see, e.g., Corollary 3.13 in [19] and Theorem 3 below), and thus all have the same cardinality.\par
 
% END S2.p5
 
% BEGIN Thmtheorem2
 
\par\medskip\noindent\textbf{Definition 2.}\par
 
% BEGIN Thmtheorem2.p1
 
Let $\Gamma$ be a solvable $E$-unification problem and $X$ a set of variables such that $\operatorname{Var}(\Gamma)\subseteq X\subseteq V$. Then the unification type of $\Gamma$ w.r.t. $\leq_{E}^{X}$ is\par
  
\par\noindent  $\blacksquare$ 
% BEGIN S2.I1.i1.p1
 
unitary if $\Gamma$ has a minimal complete set of $E$-unifiers of cardinality one w.r.t. $\leq_{E}^{X}$, whose single element is then called most general $E$-unifier (mgu),\par
 
% END S2.I1.i1.p1
 
 
\par\noindent  $\blacksquare$ 
% BEGIN S2.I1.i2.p1
 
finitary if $\Gamma$ has a finite minimal complete set of $E$-unifiers of cardinality greater than one w.r.t. $\leq_{E}^{X}$,\par
 
% END S2.I1.i2.p1
 
 
\par\noindent  $\blacksquare$ 
% BEGIN S2.I1.i3.p1
 
infinitary if $\Gamma$ has an infinite minimal complete set of $E$-unifiers w.r.t. $\leq_{E}^{X}$,\par
 
% END S2.I1.i3.p1
 
 
\par\noindent  $\blacksquare$ 
% BEGIN S2.I1.i4.p1
 
zero if $\Gamma$ does not have a minimal complete set of $E$-unifiers w.r.t. $\leq_{E}^{X}$, i.e., every complete set is redundant in the sense that it must contain two distinct elements that are comparable w.r.t. $\leq_{E}^{X}$.\par
 
% END S2.I1.i4.p1
 
  
% END Thmtheorem2.p1
 
% END Thmtheorem2
 
% BEGIN S2.p6
 
Minimal complete sets of unifiers can alternatively be characterized using the following order-theoretic point of view [3, 19]. Let $\Gamma$ be an $E$-unification problem and $X\subseteq V$ a set of variables satisfying $\operatorname{Var}(\Gamma)\subseteq X$. We denote the $\sim_{E}^{X}$-equivalence class of a unifier $\sigma$ as $[\sigma]_{E}^{X}$ and the set of all equivalence classes of unifiers as $[\mathcal{U}_{E}(\Gamma)]_{E}^{X}$. The partial order $\preceq_{E}^{X}$ on $[\mathcal{U}_{E}(\Gamma)]_{E}^{X}$ induced by the instantiation preorder $\leq_{E}^{X}$ on unifiers is defined as $[\sigma]_{E}^{X}\preceq_{E}^{X}[\tau]_{E}^{X}$ if $\sigma\leq_{E}^{X}\tau$. We say that $S\subseteq[\mathcal{U}_{E}(\Gamma)]_{E}^{X}$ is complete w.r.t. $\preceq_{E}^{X}$ if every element of $[\mathcal{U}_{E}(\Gamma)]_{E}^{X}$ is above (w.r.t. $\preceq_{E}^{X}$) some element of $S$.\par
 
% END S2.p6
 
% BEGIN Thmtheorem3
 
\par\medskip\noindent\textbf{Theorem 3 ([19]).}\par
 
% BEGIN Thmtheorem3.p1
 
Let $M$ be the set of $\preceq_{E}^{X}$-minimal elements of $[\mathcal{U}_{E}(\Gamma)]_{E}^{X}$. If $\mathcal{S}$ is a minimal complete set of $E$-unifiers of $\Gamma$ w.r.t. $\leq_{E}^{X}$, then $M=\{[\sigma]_{E}^{X}\mid\sigma\in\mathcal{S}\}$. Conversely, if $M$ is complete in $[\mathcal{U}_{E}(\Gamma)]_{E}^{X}$, then any set of substitutions obtained by picking one representative for each element of $M$ is a minimal complete set of $E$-unifiers of $\Gamma$.\par
 
% END Thmtheorem3.p1
 
% END Thmtheorem3
 
% BEGIN S2.p7
 
Consequently, unification type zero corresponds to the case where the set $M$ of minimal elements is not complete, whereas the other types are determined by the cardinality of the set $M$ in case it is complete. Theorem 3.1 in [3] establishes conditions that are necessary, sufficient, or both for proving unification type zero.333Note that, in [3], the instantiation preorder is written the other way round, i.e., more general substitutions are larger.  Here, we present one of these sufficient conditions since we will use it in Section 3 to show that the unification type of the description logic $\mathcal{E\mkern-1.618muL}$ is zero w.r.t. the restricted instantiation preorder. It was also employed in the first paper showing unification type zero for an equational theory [29, 30].\par
 
% END S2.p7
 
% BEGIN Thmtheorem4
 
\par\medskip\noindent\textbf{Lemma 4 ([3]).}\par
 
% BEGIN Thmtheorem4.p1
 
If there exists a strictly decreasing chain $\sigma_{1}>_{E}^{X}\sigma_{2}>_{E}^{X}\sigma_{3}>_{E}^{X}\cdots$ such that the set $\mathcal{S}=\{\sigma_{1},\sigma_{2},\sigma_{3},\ldots\}$ is a complete set of $E$-unifiers of $\Gamma$ w.r.t. $\leq_{E}^{X}$, then $\Gamma$ has type zero w.r.t. $\leq_{E}^{X}$.\par
 
% END Thmtheorem4.p1
 
% END Thmtheorem4
 
% BEGIN S2.SS0.SSS0.P0.SPx1
 
\par\medskip\noindent\textbf{Proof.}\par
 
% BEGIN S2.SS0.SSS0.P0.SPx1.p1
 
Suppose there exists a chain $\sigma_{1}>_{E}^{X}\sigma_{2}>_{E}^{X}\sigma_{3}>_{E}^{X}\cdots$ satisfying the conditions stated in the lemma. This chain satisfies the following properties:\par
  
\par\noindent  1. 
% BEGIN S2.I2.i1.p1
 
It has no lower bound in $\mathcal{U}_{E}(\Gamma)$, i.e., there is no $E$-unifier $\tau$ of $\Gamma$ such that $\sigma_{i}\geq_{E}^{X}\tau$ for all $i\geq 1$. To see why this is true, suppose such a unifier $\tau$ does exist. Since $\mathcal{S}$ is complete, there is $j\geq 1$ such that $\tau\geq_{E}^{X}\sigma_{j}>_{E}^{X}\sigma_{j+1}\geq_{E}^{X}\tau$. Transitivity of $\geq_{E}^{X}$ yields that $\sigma_{j+1}\geq_{E}^{X}\sigma_{j}$, but this contradicts $\sigma_{j}>_{E}^{X}\sigma_{j+1}$.\par
 
% END S2.I2.i1.p1
 
 
\par\noindent  2. 
% BEGIN S2.I2.i2.p1
 
For all $i\geq 1$, if there is $\tau\in\mathcal{U}_{E}(\Gamma)$ such that $\sigma_{i}\geq_{E}^{X}\tau$, then there exists $\tau^{\prime}\in\mathcal{U}_{E}(\Gamma)$ such that $\tau\geq_{E}^{X}\tau^{\prime}$ and $\sigma_{i+1}\geq_{E}^{X}\tau^{\prime}$. To show this, assume that $\sigma_{i}\geq_{E}^{X}\tau$. The completeness of $\mathcal{S}$ yields $j\geq 1$ such that $\sigma_{i}\geq_{E}^{X}\tau\geq_{E}^{X}\sigma_{j}$. This implies that $\sigma_{i}\geq_{E}^{X}\sigma_{j}$ because $\geq_{E}^{X}$ is a transitive relation. Hence, since the chain is strictly decreasing, it follows that $i\leq j$, and thus $\sigma_{i+1}\geq_{E}^{X}\sigma_{j+1}$ and $\tau\geq_{E}^{X}\sigma_{j+1}$. Consequently, we can set $\tau^{\prime}:=\sigma_{j+1}$.\par
 
% END S2.I2.i2.p1
 
  
Now, let $M$ be the set of $\preceq_{E}^{X}$-minimal elements of $[\mathcal{U}_{E}(\Gamma)]_{E}^{X}$. To prove that $\Gamma$ has type zero, it suffices to show that $M$ is not complete. Assume to the contrary that $M$ is complete. Then there exists $[\tau]^{X}_{E}\in M$ such that $[\tau]^{X}_{E}\preceq_{E}^{X}[\sigma_{1}]^{X}_{E}$, and thus $\sigma_{1}\geq_{E}^{X}\tau$. Since our chain does not have a lower bound, there must be an $i\geq 1$ such that $\sigma_{i}\geq_{E}^{X}\tau$ and $\sigma_{i+1}\not\geq_{E}^{X}\tau$. Using the second property shown for our chain, there exists a unifier $\tau^{\prime}$ such that $\tau\geq_{E}^{X}\tau^{\prime}$ and $\sigma_{i+1}\geq_{E}^{X}\tau^{\prime}$. Minimality of $[\tau]^{X}_{E}$ implies that $\tau\sim_{E}^{X}\tau^{\prime}$, and thus $\sigma_{i+1}\geq_{E}^{X}\tau^{\prime}\geq_{E}^{X}\tau$, which yields a contradiction. Thus, we have shown that $M$ cannot be complete, which implies that $\Gamma$ has type zero. $\hfill\blacktriangleleft$\par
 
% END S2.SS0.SSS0.P0.SPx1.p1
 
% BEGIN S2.SS0.SSS0.P0.SPx1.p2
 
As usual, we order unification types w.r.t. how bad they are (larger is worse) by setting\par
 
zero $>$ infinitary $>$ finitary $>$ unitary.\par
 
The unification type w.r.t. $\leq_{E}^{X}$ of an equational theory $E$ is the worst type of any solvable $E$-unification problem w.r.t. $\leq_{E}^{X}$. The unrestricted unification type of $\Gamma$ ($E$) is the one w.r.t. $\leq_{E}^{V}$ and the restricted unification type of $\Gamma$ ($E$) is the one w.r.t. $\leq_{E}^{X}$ for $X=\operatorname{Var}(\Gamma)$.\par
 
% END S2.SS0.SSS0.P0.SPx1.p2
 
% BEGIN S2.SS0.SSS0.P0.SPx1.p3
 
The results on the unification type of equational theories in the literature are usually shown for the restricted case. As we will demonstrate in this paper, it may indeed make a considerable difference for the unification type which instantiation preorder is employed in its definition. This is however not the case in the following situation.\par
 
% END S2.SS0.SSS0.P0.SPx1.p3
 
% BEGIN Thmtheorem5
 
\par\medskip\noindent\textbf{Lemma 5.}\par
 
% BEGIN Thmtheorem5.p1
 
Let $E$ be an equational theory, $\Gamma$ an $E$-unification problem, $X_{0}:=\operatorname{Var}(\Gamma)$, and $X\subseteq V$ a set of variables such that $X_{0}\subseteq X$ and $V\setminus X$ is infinite. If $\Gamma$ has a minimal complete set of $E$-unifiers w.r.t. $\leq_{E}^{X}$, then it has a minimal complete set of $E$-unifiers w.r.t. $\leq_{E}^{X_{0}}$ of the same cardinality, and vice versa.\par
 
% END Thmtheorem5.p1
 
% END Thmtheorem5
 
% END S2.SS0.SSS0.P0.SPx1
 
% BEGIN S2.SS0.SSS0.P0.SPx2
 
\par\medskip\noindent\textbf{Proof.}\par
 
% BEGIN S2.SS0.SSS0.P0.SPx2.p1
 
Let $\mathcal{S}$ be a minimal complete set of $E$-unifiers of $\Gamma$ w.r.t. $\leq_{E}^{X}$. For every unifier $\theta$ in $\mathcal{S}$ we can rename the variables in $\operatorname{VRan}(\theta)$ such that they do not belong to $X$ by applying an appropriate permutation $\pi$ that maps the variables of $\operatorname{VRan}(\theta)$ bijectively to a set of variables $Y$ with $|Y|=|\operatorname{VRan}(\theta)|$ and $Y\cap(X\cup\operatorname{VRan}(\theta))=\emptyset$. Such a finite set $Y$ of variables exists since $\operatorname{VRan}(\theta)$ is finite and $V\setminus X$ is infinite. For a given bijection $p:\operatorname{VRan}(\theta)\rightarrow Y$, we can define $\pi$ as follows: $\pi(z):=p(z)$ for all $z\in\operatorname{VRan}(\theta)$, $\pi(y):=p^{-1}(y)$ for all $y\in Y$, and $\pi(x)=x$ for all other variables. Note that $\pi$ is a substitution since its domain is $\operatorname{VRan}(\theta)\cup Y$, which is finite. To show that $\pi$ really is a permutation, i.e., a bijective mapping from $V$ to $V$, it is sufficient to show that it is an injective mapping from $V$ to $V$ (see Lemma 2.6 in [28]). This is an immediate consequence of the facts that $p$ and $p^{-1}$ are bijections and $Y$ and $\operatorname{VRan}(\theta)$ are disjoint.\par
 
% END S2.SS0.SSS0.P0.SPx2.p1
 
% BEGIN S2.SS0.SSS0.P0.SPx2.p2
 
The substitution $\pi\theta$ is equivalent to $\theta$ w.r.t. the equivalence relation $\sim_{\emptyset}^{V}$ induced by $\leq_{\emptyset}^{V}$, and thus also w.r.t. $\sim_{E}^{X}$. In fact, $\pi^{-1}\pi\theta=\theta$. If we restrict the domain of this substitution to $X$, then the resulting substitution $\pi\theta|_{X}$ is still equivalent to $\theta$ w.r.t. $\sim_{E}^{X}$ and also satisfies $\operatorname{VRan}(\pi\theta|_{X})\cap X=\emptyset$. To see the latter, consider a variable $x\in\operatorname{Dom}(\pi\theta|_{X})$. First assume that $x\in\operatorname{Dom}(\theta)\cap X$. Then all the variables in $\theta(x)$ belong to $\operatorname{VRan}(\theta)$, and thus all the variables in $\pi\theta|_{X}(x)$ belong to $Y$, which is disjoint with $X$. If $x\in X\setminus\operatorname{Dom}(\theta)$, then $\pi\theta|_{X}(x)=\pi(x)\neq x$. Since $\operatorname{Dom}(\pi)=\operatorname{VRan}(\theta)\cup Y$ and $Y$ is disjoint with $X$, this implies $x\in\operatorname{VRan}(\theta)$, and thus $\pi(x)\in Y$.\par
 
% END S2.SS0.SSS0.P0.SPx2.p2
 
% BEGIN S2.SS0.SSS0.P0.SPx2.p3
 
Let $\mathcal{S}^{\prime}$ be the set of substitutions obtained from $\mathcal{S}$ by applying this renaming and domain restriction process to every element of $\mathcal{S}$. Due to the $\sim_{E}^{X}$-equivalence of the elements of $\mathcal{S}$ with their modified variants in $\mathcal{S}^{\prime}$ and the fact that $X$ contains all the variables occurring in $\Gamma$, it is easy to see that $\mathcal{S}^{\prime}$ is also a minimal complete set of $E$-unifiers of $\Gamma$ w.r.t. $\leq_{E}^{X}$. In a second step, we restrict the domains of the elements of $\mathcal{S}^{\prime}$ to $X_{0}$. Given an element $\theta$ of $\mathcal{S}^{\prime}$, we define $\theta|_{X_{0}}$ to be the substitution that coincides with $\theta$ on $X_{0}$ and maps each variable $x\not\in X_{0}$ to $x$. Since $\theta$ is an $E$-unifier of $\Gamma$ and this unification problem contains only variables from $X_{0}$, the substitution $\theta|_{X_{0}}$ is clearly also an $E$-unifier of $\Gamma$. We claim that $\theta|_{X_{0}}\leq_{E}^{X}\theta$. To prove this, we define the substitution $\lambda$ by setting $\lambda(x):=\theta(x)$ if $x\in X\setminus X_{0}$ and $\lambda(x):=x$ for all other variables. We now show that $\lambda\theta|_{X_{0}}\approx_{E}^{X}\theta$. If $x\in X_{0}\cap\operatorname{Dom}(\theta)$, then $\lambda(\theta|_{X_{0}}(x))=\lambda(\theta(x))=\theta(x)$. The first identity holds by the definition of $\theta|_{X_{0}}$ and the second since the variables occurring in $\theta(x)$ are elements of $\operatorname{VRan}(\theta)$, and thus do not belong to $X$. If $x\in X_{0}\setminus\operatorname{Dom}(\theta)$, then $\lambda(\theta|_{X_{0}}(x))=\lambda(x)=x=\theta(x)$. Finally, if $x\in X\setminus X_{0}$, then $\lambda(\theta|_{X_{0}}(x))=\lambda(x)=\theta(x)$. Again, the first identity holds by the definition of $\theta|_{X_{0}}$ and the second by the definition of $\lambda$. Summing up, we have shown that the following holds for every element $\theta$ of $\mathcal{S}^{\prime}$: $\theta|_{X_{0}}$ is an $E$-unifier of $\Gamma$ and $\theta|_{X_{0}}\leq_{E}^{X}\theta$. Since $\mathcal{S}^{\prime}$ is a minimal complete set of $E$-unifiers of $\Gamma$ w.r.t. $\leq_{E}^{X}$, its elements are minimal w.r.t. $\leq_{E}^{X}$, which yields $\theta|_{X_{0}}\sim_{E}^{X}\theta$. Consequently, we know that the set $\mathcal{S}^{\prime}|_{X_{0}}:=\{\theta|_{X_{0}}\mid\theta\in\mathcal{S}^{%
\prime}\}$ is also a minimal complete set of $E$-unifiers of $\Gamma$ w.r.t. $\leq_{E}^{X}$.\par
 
% END S2.SS0.SSS0.P0.SPx2.p3
 
% BEGIN S2.SS0.SSS0.P0.SPx2.p4
 
We claim that the same is true w.r.t. the smaller set of variables $X_{0}$, i.e., that $\mathcal{S}^{\prime}|_{X_{0}}$ is also minimal and complete w.r.t. $\leq_{E}^{X_{0}}$. Completeness trivially follows from the fact that ${\leq_{E}^{X}}\subseteq{\leq_{E}^{X_{0}}}$. To prove minimality, assume that $\theta$ and $\tau$ are two distinct elements of $\mathcal{S}^{\prime}|_{X_{0}}$. Then these two substitutions are not comparable w.r.t. $\leq_{E}^{X}$. Assume that they are comparable w.r.t. $\leq_{E}^{X_{0}}$, i.e., there is a substitution $\lambda$ such that $\lambda(\theta(x))\approx_{E}\tau(x)$ holds for all $x\in X_{0}$. By our construction of $\mathcal{S}^{\prime}|_{X_{0}}$, we know that $\operatorname{Dom}(\theta)\subseteq X_{0}$, $\operatorname{Dom}(\tau)\subseteq X_{0}$, and $\operatorname{VRan}(\theta)\cap X=\emptyset$. If $x$ is a variable in $X\setminus X_{0}$, then $\lambda(\theta(x))=\lambda(x)$ and $\tau(x)=x$. Thus, if we modify $\lambda$ to $\lambda^{\prime}$ such that $\lambda^{\prime}(x)=x$ holds for all $x\in X\setminus X_{0}$, then $\lambda^{\prime}(\theta(x))=\tau(x)$ holds for all $x\in X\setminus X_{0}$. This modification has no effect on the variables $x\in X_{0}$. In fact, let $x$ be such a variable. If $x\in\operatorname{Dom}(\theta)$, then $\theta(x)$ does not contain any variable from $X$, and thus $\lambda^{\prime}(\theta(x))=\lambda(\theta(x))\approx_{E}\tau(x)$ holds. If $x\notin\operatorname{Dom}(\theta)$, then $\theta(x)=x$, and thus again $\lambda^{\prime}(\theta(x))=\lambda^{\prime}(x)=\lambda(x)=\lambda(\theta(x))%
\approx_{E}\tau(x)$ since $\lambda^{\prime}$ coincides with $\lambda$ on the variables in $X_{0}$. Summing up, we have shown that the assumption $\theta\leq_{E}^{X_{0}}\tau$ implies $\theta\leq_{E}^{X}\tau$, which contradicts the fact that $\mathcal{S}^{\prime}|_{X_{0}}$ is minimal w.r.t. $\leq_{E}^{X}$. Consequently, $\mathcal{S}^{\prime}|_{X_{0}}$ is a minimal complete set w.r.t. $\leq_{E}^{X_{0}}$, and this set has the same cardinality as the minimal complete set $\mathcal{S}$ of $E$-unifiers of $\Gamma$ w.r.t. $\leq_{E}^{X}$ we have started with.\par
 
% END S2.SS0.SSS0.P0.SPx2.p4
 
% BEGIN S2.SS0.SSS0.P0.SPx2.p5
 
Conversely, let $\mathcal{S}$ be a minimal complete set of $E$-unifiers of $\Gamma$ w.r.t. $\leq_{E}^{X_{0}}$. By applying the same construction as in the proof of the other direction, we can assume without loss of generality that every unifier $\theta\in\mathcal{S}$ satisfies $\operatorname{Dom}(\theta)\subseteq X_{0}$ and $\operatorname{VRan}(\theta)\cap X=\emptyset$. In fact, by applying this construction to a minimal complete set $\mathcal{S}^{\prime}$ of $E$-unifiers of $\Gamma$ w.r.t. $\leq_{E}^{X_{0}}$, we obtain a new set $\mathcal{S}$ where every element $\theta^{\prime}$ of the original set $\mathcal{S}^{\prime}$ is replaced with an element $\theta$ that satisfies the above conditions and is $\sim_{E}^{X}$-equivalent to $\theta^{\prime}$, and thus also $\sim_{E}^{X_{0}}$-equivalent. Consequently, this new set $\mathcal{S}$ is also a minimal complete set of $E$-unifiers of $\Gamma$ w.r.t. $\leq_{E}^{X_{0}}$.\par
 
% END S2.SS0.SSS0.P0.SPx2.p5
 
% BEGIN S2.SS0.SSS0.P0.SPx2.p6
 
We claim that $\mathcal{S}$ is also minimal and complete w.r.t. the larger set of variables $X$. Minimality trivially follows from the fact that ${\leq_{E}^{X}}\subseteq{\leq_{E}^{X_{0}}}$. To prove completeness, assumed that $\tau$ is an $E$-unifier of $\Gamma$. Completeness of $\mathcal{S}$ w.r.t. $\leq_{E}^{X_{0}}$ yields an element $\theta$ of $\mathcal{S}$ such that $\theta\leq_{E}^{X_{0}}\tau$, i.e., there is a substitution $\lambda$ such that $\lambda\theta(x)\approx_{E}\tau(x)$ holds for all $x\in X_{0}$. Let $x$ be a variable in $X\setminus X_{0}$. Then $\lambda\theta(x)=\lambda(x)$. Thus, if we modify $\lambda$ to $\lambda^{\prime}$ such that $\lambda^{\prime}(x)=\tau(x)$ holds for all $x\in X\setminus X_{0}$, then $\lambda^{\prime}(\theta(x))=\tau(x)$ holds for all $x\in X\setminus X_{0}$. The claim that this modification has no effect for the variables $x\in X_{0}$ can be shown as in the proof of minimality of $\mathcal{S}^{\prime}|_{X_{0}}$ w.r.t. $\leq_{E}^{X_{0}}$ above. This proves that $\lambda^{\prime}\theta(x)\approx_{E}\tau(x)$ holds for all $x\in X$, and thus $\mathcal{S}$ is also complete w.r.t. $\leq_{E}^{X}$. $\hfill\blacktriangleleft$\par
 
% END S2.SS0.SSS0.P0.SPx2.p6
 
% BEGIN S2.SS0.SSS0.P0.SPx2.p7
 
The following theorem is an immediate consequence of this lemma.\par
 
% END S2.SS0.SSS0.P0.SPx2.p7
 
% BEGIN Thmtheorem6
 
\par\medskip\noindent\textbf{Theorem 6.}\par
 
% BEGIN Thmtheorem6.p1
 
Let $E$ be an equational theory, $\Gamma$ an $E$-unification problem, and $X\subseteq V$ a set of variables such that $\operatorname{Var}(\Gamma)\subseteq X$ and $V\setminus X$ is infinite. Then the restricted unification type of $\Gamma$ coincides with the unification type of $\Gamma$ w.r.t. $\leq_{E}^{X}$.\par
 
% END Thmtheorem6.p1
 
% END Thmtheorem6
 
% BEGIN S2.SS0.SSS0.P0.SPx2.p8
 
Note that the condition of the theorem is in particular satisfied if $X$ is a finite superset of $\operatorname{Var}(\Gamma)$. However, it is clearly not satisfied for $X=V$, which corresponds to the unrestricted instantiation preorder setting. We will see below that in this case the unification type can indeed depend on whether the restricted or the unrestricted instantiation preorder is used. However, there is a special case where this cannot happen.\par
 
% END S2.SS0.SSS0.P0.SPx2.p8
 
% BEGIN Thmtheorem7
 
\par\medskip\noindent\textbf{Theorem 7.}\par
 
% BEGIN Thmtheorem7.p1
 
Let $E$ be an equational theory, $\Gamma$ an $E$-unification problem, and $\mathcal{S}$ a set of $E$-unifiers of $\Gamma$ such that $\operatorname{VRan}(\sigma)\cup\operatorname{Dom}(\sigma)\subseteq%
\operatorname{Var}(\Gamma)$ holds for all $\sigma\in\mathcal{S}$. Then $\mathcal{S}$ is a minimal complete set of $E$-unifiers of $\Gamma$ w.r.t. $\leq_{E}^{V}$ iff it is a minimal complete set of $E$-unifiers of $\Gamma$ w.r.t. $\leq_{E}^{\operatorname{Var}(\Gamma)}$.\par
 
% END Thmtheorem7.p1
 
% END Thmtheorem7
 
% END S2.SS0.SSS0.P0.SPx2
 
% BEGIN S2.SS0.SSS0.P0.SPx3
 
\par\medskip\noindent\textbf{Proof.}\par
 
% BEGIN S2.SS0.SSS0.P0.SPx3.p1
 
Let $\mathcal{S}$ be a minimal complete set of $E$-unifiers of $\Gamma$ w.r.t. $\leq_{E}^{V}$. Since ${\leq_{E}^{V}}\subseteq{\leq_{E}^{\operatorname{Var}(\Gamma)}}$, this set is also complete w.r.t. $\leq_{E}^{\operatorname{Var}(\Gamma)}$. To show minimality w.r.t. $\leq_{E}^{\operatorname{Var}(\Gamma)}$, assume to the contrary that $\sigma,\theta$ are two distinct elements of $\mathcal{S}$ such that $\sigma\leq_{E}^{\operatorname{Var}(\Gamma)}\theta$, i.e., there is a substitution $\lambda$ such that $\lambda\sigma(x)\approx_{E}\theta(x)$ holds for all $x\in\operatorname{Var}(\Gamma)$. We modify $\lambda$ to $\lambda^{\prime}$ by setting $\lambda^{\prime}(x)=x$ for all variables $x\in V\setminus\operatorname{Var}(\Gamma)$. For $x\in\operatorname{Var}(\Gamma)$, we know that $\sigma(x)$ contains only variables from $\operatorname{VRan}(\sigma)\subseteq\operatorname{Var}(\Gamma)$ if $x\in\operatorname{Dom}(\sigma)$ or $\sigma(x)=x\in\operatorname{Var}(\Gamma)$. Since $\lambda$ and $\lambda^{\prime}$ coincide on $\operatorname{Var}(\Gamma)$, this yields $\lambda^{\prime}\sigma(x)=\lambda\sigma(x)\approx_{E}\theta(x)$. For $x\in V\setminus\operatorname{Var}(\Gamma)$, this variable does not belong to any of the sets $\operatorname{Dom}(\sigma)$, $\operatorname{Dom}(\theta)$, and $\operatorname{Dom}(\lambda^{\prime})$, and thus $\lambda^{\prime}\sigma(x)=\lambda^{\prime}(x)=x=\theta(x)$. Summing up, we have shown that $\sigma\leq_{E}^{V}\theta$, which contradicts our assumption that $\mathcal{S}$ is minimal w.r.t. $\leq_{E}^{V}$.\par
 
% END S2.SS0.SSS0.P0.SPx3.p1
 
% BEGIN S2.SS0.SSS0.P0.SPx3.p2
 
Conversely, assume that $\mathcal{S}$ is a minimal complete set of $E$-unifiers of $\Gamma$ w.r.t. $\leq_{E}^{\operatorname{Var}(\Gamma)}$. Since ${\leq_{E}^{V}}\subseteq{\leq_{E}^{\operatorname{Var}(\Gamma)}}$, minimality also holds w.r.t. $\leq_{E}^{V}$. To show completeness w.r.t. $\leq_{E}^{V}$, assume that $\theta$ is an $E$-unifier of $\Gamma$. Completeness of $\mathcal{S}$ w.r.t. $\leq_{E}^{\operatorname{Var}(\Gamma)}$ yields a substitution $\sigma\in\mathcal{S}$ such that $\sigma\leq_{E}^{\operatorname{Var}(\Gamma)}\theta$. Similarly to the first part of the proof, we can show that this also implies $\sigma\leq_{E}^{V}\theta$. The difference is that now $\theta$ is an arbitrary unifier, and thus $\operatorname{VRan}(\theta)\cup\operatorname{Dom}(\theta)\subseteq%
\operatorname{Var}(\Gamma)$ need not hold. Let $\lambda$ be such that $\lambda\sigma(x)\approx_{E}\theta(x)$ holds for all $x\in\operatorname{Var}(\Gamma)$. We modify $\lambda$ to $\lambda^{\prime}$ by setting $\lambda^{\prime}(x)=\theta(x)$ for all variables $x\in V\setminus\operatorname{Var}(\Gamma)$. For $x\in\operatorname{Var}(\Gamma)$, we obtain $\lambda^{\prime}\sigma(x)=\lambda\sigma(x)\approx_{E}\theta(x)$ as in the first part of the proof. For $x\in V\setminus\operatorname{Var}(\Gamma)$, this variable does not belong to $\operatorname{Dom}(\sigma)$, and thus $\lambda^{\prime}\sigma(x)=\lambda^{\prime}(x)=\theta(x)$. $\hfill\blacktriangleleft$\par
 
% END S2.SS0.SSS0.P0.SPx3.p2
 
% BEGIN S2.SS0.SSS0.P0.SPx3.p3
 
Examples of equational theories where the conditions of this theorem are always satisfied are the empty theory (unitary), the theory $\mathsf{C}$ axiomatizing commutativity of a binary function symbol (finitary), and the theory $\mathsf{A}$ axiomatizing associativity of a binary function symbol (infinitary). This is an easy consequence of the known algorithms [42, 50, 45] computing (or enumerating, in the case of $\mathsf{A}$) minimal complete sets of unifiers for these theories.\par
 
% END S2.SS0.SSS0.P0.SPx3.p3
 
% BEGIN S2.SS0.SSS0.P0.SPx3.p4
 
The following result is an easy consequence of the fact that ${\leq_{E}^{V}}\subseteq{\leq_{E}^{\operatorname{Var}(\Gamma)}}$.\par
 
% END S2.SS0.SSS0.P0.SPx3.p4
 
% BEGIN Thmtheorem8
 
\par\medskip\noindent\textbf{Theorem 8.}\par
 
% BEGIN Thmtheorem8.p1
 
Let $E$ be an equational theory. If $E$ has unrestricted unification type unitary (finitary), then it has restricted unification type unitary (finitary or unitary).\par
 
% END Thmtheorem8.p1
 
% END Thmtheorem8
 
% END S2.SS0.SSS0.P0.SPx3
 
% BEGIN S2.SS0.SSS0.P0.SPx4
 
\par\medskip\noindent\textbf{Proof.}\par
 
% BEGIN S2.SS0.SSS0.P0.SPx4.p1
 
A given finite minimal complete set $\mathcal{S}$ of $E$-unifiers of $\Gamma$ w.r.t. $\leq_{E}^{V}$ is also complete w.r.t. $\leq_{E}^{\operatorname{Var}(\Gamma)}$. If it is not minimal (which can only happen in the finitary case), then one can make it minimal by removing redundant element (i.e., elements $\theta$ such that $\mathcal{S}$ contains an element $\sigma<_{E}^{\operatorname{Var}(\Gamma)}\theta$) without destroying completeness. $\hfill\blacktriangleleft$\par
 
% END S2.SS0.SSS0.P0.SPx4.p1
 
% END S2.SS0.SSS0.P0.SPx4
 
% END S2

% BEGIN S3
 
\par\medskip\noindent\textbf{3 The unification type of $\mathcal{E\mkern-1.618muL}$}\par
 
% BEGIN S3.p1
 
Unification was introduced in description logics as a tool for detecting redundancies in large knowledge bases [14]. The description logic $\mathcal{E\mkern-1.618muL}$ [6] has drawn considerable attention since its standard reasoning problems can be solved in polynomial time while the logic is still expressive enough for formalizing bio-medical ontologies [11]. In [12, 13] it was shown that, in the setting of unification with constants, $\mathcal{E\mkern-1.618muL}$ has unification type zero w.r.t. the restricted instantiation preorder. In the following, we will analyze the example used in [13] to prove this result in more detail, and explain why it does not work for the unrestricted instantiation preorder. Then we show that the unification type of $\mathcal{E\mkern-1.618muL}$ in the unrestricted setting is actually infinitary, not just for this example, but in general.\par
 
% END S3.p1
 
% BEGIN S3.p2
 
But first, we briefly introduce $\mathcal{E\mkern-1.618muL}$ and recall why unification in $\mathcal{E\mkern-1.618muL}$ can be seen as unification modulo an equational theory (see [6, 13, 15] for more detailed descriptions). Given sets of concept names (unary predicates) and role names (binary predicates), $\mathcal{E\mkern-1.618muL}$ concept descriptions (or simply concepts) are built from concept names using the concept constructors top concept ($\top$), conjunction ($C\sqcap D$), and existential restriction ($\exists r.C$). In the model-theoretic semantics of $\mathcal{E\mkern-1.618muL}$, a given interpretation $\mathcal{I}$ assigns sets $C^{\mathcal{I}}$ to concept descriptions $C$ according to the semantics of the constructors. To be more precise, an interpretation $\mathcal{I}=(\Delta^{\mathcal{I}},\cdot^{\mathcal{I}})$ consists of a non-empty interpretation domain and an extension function that assigns subsets of this domain to concept names and binary relations on the domain to role names. This interpretation function is extended to concept descriptions as follows:\par
    \[\top^{\mathcal{I}}:=\Delta^{\mathcal{I}},\ \ (C\sqcap D)^{\mathcal{I}}:=C^{%
\mathcal{I}}\cap D^{\mathcal{I}},\ \ (\exists r.C)^{\mathcal{I}}:=\{d\in\Delta%
^{\mathcal{I}}\mid\exists e\in C^{\mathcal{I}}\ \mbox{such that}\ (d,e)\in r^{%
\mathcal{I}}\}.\]    
The concept description $C$ is subsumed by the concept description $D$ (written $C\sqsubseteq D$) if $C^{\mathcal{I}}\subseteq D^{\mathcal{I}}$ holds for all interpretations $\mathcal{I}$, and $C$ and $D$ are equivalent (written $C\equiv D$) if they subsume each other.\par
 
% END S3.p2
 
% BEGIN S3.p3
 
In [13], the notion of a reduced $\mathcal{E\mkern-1.618muL}$ concept is employed to derive characterizations of equivalence and subsumption of $\mathcal{E\mkern-1.618muL}$ concepts. Here, we recall this notion and the characterizations since they are used later on to explain why $\mathcal{E\mkern-1.618muL}$ has unification type zero w.r.t. the restricted instantiation preorder. Küsters [39] introduces the following rules for reducing $\mathcal{E\mkern-1.618muL}$ concept descriptions:\par
    \[\begin{array}[]{l@{\ \ \ \ }l}C\sqcap\top\rightarrow C&\mbox{for all $\mathcal%
{EL}$ concept descriptions $C$},\\[3.00003pt]
A\sqcap A\rightarrow A&\mbox{for all concept names $A$},\\[3.00003pt]
\exists r.C\sqcap\exists r.D\rightarrow\exists r.C&\mbox{for all $\mathcal{EL}%
$ concept descriptions $C,D$
with $C\sqsubseteq D$}.\end{array}\]    
A reduced form of a given $\mathcal{E\mkern-1.618muL}$ concept description $C$ is then obtained from $C$ by applying these rules exhaustively modulo associativity and commutativity of $\sqcap$. The following theorem is stated in [13] as an easy consequence of Corollary 6.3.1 on page 181 of [39].\par
 
% END S3.p3
 
% BEGIN Thmtheorem9
 
\par\medskip\noindent\textbf{Theorem 9 (Theorem 3.1 of [13]).}\par
 
% BEGIN Thmtheorem9.p1
 
Let $C,D$ be $\mathcal{EL}$ concept descriptions, and $\widehat{C},\widehat{D}$ reduced forms of $C,D$, respectively. Then $C\equiv D$ iff $\widehat{C}$ is identical to $\widehat{D}$ up to associativity and commutativity of $\sqcap$.\par
 
% END Thmtheorem9.p1
 
% END Thmtheorem9
 
% BEGIN S3.p4
 
This theorem is used in [13] to derive a recursive characterization of subsumption in $\mathcal{EL}$.\par
 
% END S3.p4
 
% BEGIN Thmtheorem10
 
\par\medskip\noindent\textbf{Corollary 10 (Corollary 3.2 of [13]).}\par
 
% BEGIN Thmtheorem10.p1
 
Let $C=A_{1}\sqcap\ldots\sqcap A_{k}\sqcap\exists r_{1}.C_{1}\sqcap\ldots\sqcap%
\exists r_{m}.C_{m}$ and $D=B_{1}\sqcap\ldots\sqcap B_{\ell}\sqcap\exists s_{1}.D_{1}\sqcap\ldots\sqcap%
\exists s_{n}.D_{n}$, where $A_{1},\ldots,A_{k},B_{1},\ldots,B_{\ell}$ are concept names. Then $C\sqsubseteq D$ iff $\{B_{1},\ldots,B_{\ell}\}\subseteq\{A_{1},\ldots,A_{k}\}$ and for every $j,1\leq j\leq n$, there exists an $i,1\leq i\leq m$, such that $r_{i}=s_{j}$ and $C_{i}\sqsubseteq D_{j}$.\par
 
% END Thmtheorem10.p1
 
% END Thmtheorem10
 
% BEGIN S3.p5
 
From this corollary, the following lemma is then derived in [13].\par
 
% END S3.p5
 
% BEGIN Thmtheorem11
 
\par\medskip\noindent\textbf{Lemma 11 (Lemma 3.3 of [13]).}\par
 
% BEGIN Thmtheorem11.p1
 
If $C,D$ are reduced $\mathcal{EL}$ concept descriptions such that $\exists r.D\sqsubseteq C$, then $C$ is either $\top$, or of the form $C=\exists r.C_{1}\sqcap\ldots\sqcap\exists r.C_{n}$ where $n\geq 1$; $C_{1},\ldots,C_{n}$ are reduced and pairwise incomparable w.r.t. subsumption; and $D\sqsubseteq C_{1},\ldots,D\sqsubseteq C_{n}$. Conversely, if $C,D$ are $\mathcal{EL}$ concept descriptions such that $C=\exists r.C_{1}\sqcap\ldots\sqcap\exists r.C_{n}$ and $D\sqsubseteq C_{1},\ldots,D\sqsubseteq C_{n}$, then $\exists r.D\sqsubseteq C$.\par
 
% END Thmtheorem11.p1
 
% END Thmtheorem11
 
% BEGIN S3.p6
 
Equivalence of $\mathcal{E\mkern-1.618muL}$ concept descriptions can be axiomatized by the equational theory $\mathsf{bSLmO}$ of bounded semilattices with monotone operators [54, 15, 55]. For this purpose, we view the conjunction operator $\sqcap$ as a binary function symbol, $\top$ as a constant symbol, and $\exists r.$ for a role name $r$ as a unary function symbol. The theory $\mathsf{bSLmO}$ over this signature then consists of the identities stating that $\sqcap$ is associative, commutative, and idempotent, has $\top$ as unit, and existential restrictions as monotone operators:\par
    \[x\sqcap y\approx y\sqcap x,\ (x\sqcap y)\sqcap z\approx x\sqcap(y\sqcap z),\ x%
\sqcap x\approx x,\ x\sqcap\top\approx x,\exists r.x\sqcap\exists r.(x\sqcap y%
)\approx\exists r.(x\sqcap y).\]    
% END S3.p6
 
% BEGIN S3.p7
 
For unification in $\mathcal{E\mkern-1.618muL}$ we partition the set of concept names into concept variables and concept constants. Substitutions can then replace concept variables in $\mathcal{E\mkern-1.618muL}$ concept descriptions with $\mathcal{E\mkern-1.618muL}$ concept descriptions. Unifiers are supposed to solve equations between concept descriptions with variables by making them equivalent. This corresponds to unification with constants modulo the equational theory $\mathsf{bSLmO}$. In the following, we use $\approx_{\mathcal{E\mkern-1.618muL}}$ (rather than $\approx_{\mathsf{bSLmO}}$) to denote equivalence between $\mathcal{E\mkern-1.618muL}$ concept descriptions and also employ the subscript $\mathcal{E\mkern-1.618muL}$ when writing the respective instantiation preorders.\par
 
% END S3.p7
 
% BEGIN Thmtheorem12
 
\par\medskip\noindent\textbf{Example 12.}\par
 
% BEGIN Thmtheorem12.p1
 
Consider the $\mathcal{E\mkern-1.618muL}$-unification problem $\Gamma:=\{x\sqcap\exists r.y\approx_{\mathcal{E\mkern-1.618muL}}^{?}\exists r.y\}$. Then, for every $n\geq 0$, the substitution\par
    \[\sigma_{n}:=\{x\mapsto\exists r.z_{1}\sqcap\ldots\sqcap\exists r.z_{n},\ \ y%
\mapsto z_{1}\sqcap\ldots\sqcap z_{n}\sqcap z\}\]    
for distinct variables $x,y,z_{1},\ldots,z_{n},z$ is an $\mathcal{E\mkern-1.618muL}$-unifier of $\Gamma$, where the empty conjunction is $\top$, i.e., $\sigma_{0}=\{x\mapsto\top,y\mapsto z\}$. We will show below that, w.r.t. the restricted instantiation preorder, the set $\{\sigma_{n}\mid n\geq 0\}$ is a complete set of $\mathcal{E\mkern-1.618muL}$-unifiers that constitutes a strictly decreasing chain $\sigma_{0}>_{\mathcal{E\mkern-1.618muL}}^{\{x,y\}}\sigma_{1}>_{\mathcal{E%
\mkern-1.618muL}}^{\{x,y\}}\sigma_{2}>_{\mathcal{E\mkern-1.618muL}}^{\{x,y\}}\cdots$ of more and more general unifiers. According to Lemma 4 this implies that $\Gamma$ has unification type zero, and thus $\mathcal{E\mkern-1.618muL}$ has unification type zero.\par
 
% END Thmtheorem12.p1
 
% END Thmtheorem12
 
% BEGIN Thmtheorem13
 
\par\medskip\noindent\textbf{Lemma 13.}\par
 
% BEGIN Thmtheorem13.p1
 
Let $\Gamma$ and $\sigma_{n}$ for $n\geq 0$ be defined as in Example 12. Then the set $\{\sigma_{n}\mid n\geq 0\}$ is a complete set of $\mathcal{E\mkern-1.618muL}$-unifiers of $\Gamma$ w.r.t. $\leq_{\mathcal{E\mkern-1.618muL}}^{\{x,y\}}$.\par
 
% END Thmtheorem13.p1
 
% END Thmtheorem13
 
% BEGIN S3.SS0.SSS0.P0.SPx1
 
\par\medskip\noindent\textbf{Proof.}\par
 
% BEGIN S3.SS0.SSS0.P0.SPx1.p1
 
Let $\sigma$ be an $\mathcal{E\mkern-1.618muL}$-unifier of $\Gamma$ and $C:=\sigma(y)$ and $D:=\sigma(x)$. Then $\exists r.C\sqsubseteq D$. In addition, we can assume without loss of generality that $C$ and $D$ are reduced. The characterization of subsumption in $\mathcal{E\mkern-1.618muL}$ (see Corollary 10) yields $D=\exists r.D_{1}\sqcap\ldots\sqcap\exists r.D_{n}$ for $n\geq 0$ $\mathcal{E\mkern-1.618muL}$ concept descriptions $D_{1},\ldots,D_{n}$ satisfying $C\sqsubseteq D_{1},\ldots,C\sqsubseteq D_{n}$ (see Lemma 11). Thus, if we define $\lambda:=\{z_{1}\mapsto D_{1},\ldots,z_{n}\mapsto D_{n},z\mapsto C\}$, then $\lambda\sigma_{n}\approx_{\mathcal{E\mkern-1.618muL}}^{\{x,y\}}\sigma$, which shows $\sigma_{n}\leq_{\mathcal{E\mkern-1.618muL}}^{\{x,y\}}\sigma$. $\hfill\blacktriangleleft$\par
 
% END S3.SS0.SSS0.P0.SPx1.p1
 
% BEGIN Thmtheorem14
 
\par\medskip\noindent\textbf{Lemma 14.}\par
 
% BEGIN Thmtheorem14.p1
 
If $\sigma_{n}$ for $n\geq 0$ are defined as in Example 12, then $\sigma_{0}>_{\mathcal{E\mkern-1.618muL}}^{\{x,y\}}\sigma_{1}>_{\mathcal{E%
\mkern-1.618muL}}^{\{x,y\}}\sigma_{2}>_{\mathcal{E\mkern-1.618muL}}^{\{x,y\}}\cdots$.\par
 
% END Thmtheorem14.p1
 
% END Thmtheorem14
 
% END S3.SS0.SSS0.P0.SPx1
 
% BEGIN S3.SS0.SSS0.P0.SPx2
 
\par\medskip\noindent\textbf{Proof.}\par
 
% BEGIN S3.SS0.SSS0.P0.SPx2.p1
 
First note that $\sigma_{n+1}\leq_{\mathcal{E\mkern-1.618muL}}^{\{x,y\}}\sigma_{n}$ since $\lambda\sigma_{n+1}\approx_{\mathcal{E\mkern-1.618muL}}^{\{x,y\}}\sigma_{n}$ for $\lambda:=\{z_{n+1}\mapsto\top\}$. Second, assume that $\sigma_{n}\leq_{\mathcal{E\mkern-1.618muL}}^{\{x,y\}}\sigma_{n+1}$, i.e., there is a substitution $\tau$ such that $\tau\sigma_{n}\approx_{\mathcal{E\mkern-1.618muL}}^{\{x,y\}}\sigma_{n+1}$. Then $\exists r.\tau(z_{1})\sqcap\ldots\sqcap\exists r.\tau(z_{n})\approx_{\mathcal{%
E\mkern-1.618muL}}\exists r.z_{1}\sqcap\ldots\sqcap\exists r.z_{n+1}$. Consequently, by Theorem 9, the reduced forms of these two concept descriptions would need to be equal up to associativity and commutativity of $\sqcap$. However, the reduced form of the concept description on the left-hand side has at most $n$ existential restrictions, whereas the reduced form of the concept description on the right-hand side has $n+1$ existential restrictions, which yields a contradiction. Thus, we have shown $\sigma_{n+1}<_{\mathcal{E\mkern-1.618muL}}^{\{x,y\}}\sigma_{n}$. $\hfill\blacktriangleleft$\par
 
% END S3.SS0.SSS0.P0.SPx2.p1
 
% BEGIN S3.SS0.SSS0.P0.SPx2.p2
 
The following theorem is now an immediate consequence of Lemma 4.\par
 
% END S3.SS0.SSS0.P0.SPx2.p2
 
% BEGIN Thmtheorem15
 
\par\medskip\noindent\textbf{Theorem 15.}\par
 
% BEGIN Thmtheorem15.p1
 
The restricted unification type of $\mathcal{E\mkern-1.618muL}$ is zero.\par
 
% END Thmtheorem15.p1
 
% END Thmtheorem15
 
% BEGIN S3.SS0.SSS0.P0.SPx2.p3
 
With respect to the unrestricted instantiation preorder, the instance relationship between the substitutions $\sigma_{n}$ and $\sigma_{n+1}$ no longer holds. More generally, we can show the following result.\par
 
% END S3.SS0.SSS0.P0.SPx2.p3
 
% BEGIN Thmtheorem16
 
\par\medskip\noindent\textbf{Lemma 16.}\par
 
% BEGIN Thmtheorem16.p1
 
If $\sigma_{n}$ for $n\geq 0$ are defined as in Example 12, then $\sigma_{n}\not\leq_{\mathcal{E\mkern-1.618muL}}^{V}\sigma_{m}$ for all distinct $n,m\geq 0$.\par
 
% END Thmtheorem16.p1
 
% END Thmtheorem16
 
% END S3.SS0.SSS0.P0.SPx2
 
% BEGIN S3.SS0.SSS0.P0.SPx3
 
\par\medskip\noindent\textbf{Proof.}\par
 
% BEGIN S3.SS0.SSS0.P0.SPx3.p1
 
Let $n<m$. Then we know from Lemma 14 that $\sigma_{n}\not\leq_{\mathcal{E\mkern-1.618muL}}^{\{x,y\}}\sigma_{m}$, which implies $\sigma_{n}\not\leq_{\mathcal{E\mkern-1.618muL}}^{V}\sigma_{m}$ since ${\leq_{\mathcal{E\mkern-1.618muL}}^{V}}\subseteq{\leq_{\mathcal{E\mkern-1.618%
muL}}^{\{x,y\}}}$.\par
 
% END S3.SS0.SSS0.P0.SPx3.p1
 
% BEGIN S3.SS0.SSS0.P0.SPx3.p2
 
Assume that $n>m$ and that there is a substitution $\lambda$ such that $\lambda\sigma_{n}\approx_{\mathcal{E\mkern-1.618muL}}\sigma_{m}$. Then $\lambda\sigma_{n}(x)\approx_{\mathcal{E\mkern-1.618muL}}\sigma_{m}(x)$, and since $\sigma_{m}(x)$ does not contain $z_{n}$, but $\sigma_{n}(x)$ does, $\lambda$ must replace $z_{n}$ by a concept description not containing $z_{n}$, and thus in particular $\lambda(z_{n})\not\approx_{\mathcal{E\mkern-1.618muL}}z_{n}$. But then $\lambda\sigma_{n}(z_{n})=\lambda(z_{n})\not\approx_{\mathcal{E\mkern-1.618muL}%
}z_{n}$ and $\sigma_{m}(z_{n})=z_{n}$ yield a contradiction to $\lambda\sigma_{n}\approx_{\mathcal{E\mkern-1.618muL}}\sigma_{m}$. $\hfill\blacktriangleleft$\par
 
% END S3.SS0.SSS0.P0.SPx3.p2
 
% BEGIN S3.SS0.SSS0.P0.SPx3.p3
 
Thus, w.r.t. the unrestricted instantiation preorder, the set $\{\sigma_{n}\mid n\geq 0\}$ consists of unifiers that are incomparable. However, this set is also no longer complete.\par
 
% END S3.SS0.SSS0.P0.SPx3.p3
 
% BEGIN Thmtheorem17
 
\par\medskip\noindent\textbf{Lemma 17.}\par
 
% BEGIN Thmtheorem17.p1
 
Let $\Gamma$ and $\sigma_{n}$ for $n\geq 0$ be defined as in Example 12. The $\mathcal{E\mkern-1.618muL}$-unifier $\sigma:=\{x\mapsto\exists r.A,y\mapsto A\}$ of $\Gamma$ is not an instance of any of the substitutions $\sigma_{n}$ for $n\geq 0$ if we use the unrestricted instantiation preorder.\par
 
% END Thmtheorem17.p1
 
% END Thmtheorem17
 
% END S3.SS0.SSS0.P0.SPx3
 
% BEGIN S3.SS0.SSS0.P0.SPx4
 
\par\medskip\noindent\textbf{Proof.}\par
 
% BEGIN S3.SS0.SSS0.P0.SPx4.p1
 
First, note that $\sigma_{0}\leq_{\mathcal{E\mkern-1.618muL}}^{V}\sigma$ is not possible since $\sigma_{0}(x)=\top$, and thus $\lambda\sigma_{0}(x)=\top$ holds for all substitutions $\lambda$, whereas $\sigma(x)=\exists r.A$. If $n\geq 1$ and $\lambda\sigma_{n}\approx_{\mathcal{E\mkern-1.618muL}}\sigma$, then $\lambda$ must replace all variables $z_{i}$ occurring in $\sigma_{n}(x)$ with $A$, and thus $\lambda\sigma_{n}(z_{i})=\lambda(z_{i})\approx_{\mathcal{E\mkern-1.618muL}}A$. However, $\sigma(z_{i})=z_{i}$, which yields a contradiction to $\lambda\sigma_{n}\approx_{\mathcal{E\mkern-1.618muL}}\sigma$. $\hfill\blacktriangleleft$\par
 
% END S3.SS0.SSS0.P0.SPx4.p1
 
% BEGIN S3.SS0.SSS0.P0.SPx4.p2
 
Summing up, we have seen that the $\mathcal{E\mkern-1.618muL}$-unification problem $\Gamma$ of Example 12 has unification type zero for the restricted instantiation preorder, but the proof that we have used to show this result does not work if we employ the unrestricted instantiation preorder instead.\par
 
% END S3.SS0.SSS0.P0.SPx4.p2
 
% BEGIN S3.SS0.SSS0.P0.SPx4.p3
 
In the following, we prove the general result that the unrestricted unification type of $\mathcal{E\mkern-1.618muL}$ is not zero. Note that, by Theorem 8, this type cannot be unitary or finitary, and thus must be infinitary. The main idea underlying our proof of this result is to show that, up to the equivalence relation $\sim_{\mathcal{E\mkern-1.618muL}}^{V}$, every substitution has only finitely many more general substitutions w.r.t. $\leq_{\mathcal{E\mkern-1.618muL}}^{V}$. The most challenging technical result needed to prove this is the following lemma.\par
 
% END S3.SS0.SSS0.P0.SPx4.p3
 
% BEGIN Thmtheorem18
 
\par\medskip\noindent\textbf{Lemma 18.}\par
 
% BEGIN Thmtheorem18.p1
 
Let $\sigma,\theta$ be substitutions such that $\sigma\leq_{\mathcal{E\mkern-1.618muL}}^{V}\theta$. Then there is a substitution $\sigma^{\prime}$ such that $\sigma\sim_{\mathcal{E\mkern-1.618muL}}^{V}\sigma^{\prime}$ and $\operatorname{Dom}(\sigma^{\prime})\subseteq\operatorname{Dom}(\theta)$.\par
 
% END Thmtheorem18.p1
 
% END Thmtheorem18
 
% END S3.SS0.SSS0.P0.SPx4
 
% BEGIN S3.SS0.SSS0.P0.SPx5
 
\par\medskip\noindent\textbf{Proof.}\par
 
% BEGIN S3.SS0.SSS0.P0.SPx5.p1
 
First note that we can assume without loss of generality that $\operatorname{VRan}(\sigma)\cap\operatorname{Dom}(\theta)=\emptyset$. Otherwise, we can apply a permutation to $\sigma$ that renames the variables in $\operatorname{VRan}(\sigma)$ appropriately, which yields a substitution that is $\sim_{\mathcal{E\mkern-1.618muL}}^{V}$-equivalent to $\sigma$ and satisfies the required disjointness condition (see the proof of Lemma 5 for how such a permutation can be obtained).\par
 
% END S3.SS0.SSS0.P0.SPx5.p1
 
% BEGIN S3.SS0.SSS0.P0.SPx5.p2
 
Since $\sigma\leq_{\mathcal{E\mkern-1.618muL}}^{V}\theta$, we know that there is a substitution $\lambda$ such that $\lambda\sigma\approx_{\mathcal{E\mkern-1.618muL}}\theta$. Consider the set $\mathcal{X}$ of all variables $x$ such that $x\in\operatorname{Dom}(\sigma)\setminus\operatorname{Dom}(\theta)$. Then $\lambda\sigma(x)\approx_{\mathcal{E\mkern-1.618muL}}\theta(x)=x$ holds for all $x\in\mathcal{X}$. Consequently, the top-level conjunction of $\sigma(x)$ cannot contain concept constants (i.e., constants) or existential restrictions (i.e., terms starting with a function application). This means that $\sigma(x)$ is a conjunction of variables:\par
    \[\sigma(x)=z_{1}^{x}\sqcap\ldots\sqcap z_{n_{x}}^{x},\]    
where we assume (without loss of generality) that all the variables $z_{i}^{x}$ for a fixed $x$ are pairwise distinct. To obtain $\lambda\sigma(x)\approx_{\mathcal{E\mkern-1.618muL}}x$, the substitution $\lambda$ must thus assign (modulo $\approx_{\mathcal{E\mkern-1.618muL}}$) $x$ or $\top$ to the variables $z_{i}^{x}$ for all $i,1\leq i\leq n_{x}$, and $x$ to at least one of these variables. Let us assume without loss of generality that $\lambda$ assigns $x$ to $z_{1}^{x}$ for all $x\in\mathcal{X}$. This implies that $z_{1}^{x}\neq z_{1}^{y}$ for different elements $x,y$ of $\mathcal{X}$.\par
 
% END S3.SS0.SSS0.P0.SPx5.p2
 
% BEGIN S3.SS0.SSS0.P0.SPx5.p3
 
Since $\operatorname{VRan}(\sigma)\cap\operatorname{Dom}(\theta)=\emptyset$, we know that $z_{i}^{x}\not\in\operatorname{Dom}(\theta)$ for all $x\in\mathcal{X}$ and $1\leq i\leq n_{x}$. We claim that $z_{i}^{x}\in\operatorname{Dom}(\sigma)$ also holds, and thus $z_{i}^{x}\in\mathcal{X}$. Otherwise, $\lambda\sigma(z_{i}^{x})=\lambda(z_{i}^{x})\in\{x,\top\}$. However, $\lambda\sigma(z_{i}^{x})\approx_{\mathcal{E\mkern-1.618muL}}\theta(z_{i}^{x})=%
z_{i}^{x}$. This yields a contradiction unless $z_{i}^{x}=x$. But then we also have $z_{i}^{x}\in\operatorname{Dom}(\sigma)$ since $x\in\operatorname{Dom}(\sigma)$. We have thus shown that all the variables $z_{i}^{x}$ for $x\in\mathcal{X}$ also belong to $\mathcal{X}=\operatorname{Dom}(\sigma)\setminus\operatorname{Dom}(\theta)$. Since $z_{1}^{x}\neq z_{1}^{y}$ for different elements $x,y$ of $\mathcal{X}$, this implies that the $z$-variables with index $1$ already “use up” all of $\mathcal{X}$, and thus $n_{x}=1$ holds for all $x\in\mathcal{X}$.\par
 
% END S3.SS0.SSS0.P0.SPx5.p3
 
% BEGIN S3.SS0.SSS0.P0.SPx5.p4
 
Consequently, $\sigma$ is a $W$-renaming for $W=\operatorname{Dom}(\sigma)\setminus\operatorname{Dom}(\theta)$, where according to Definition 2.11 in [28] a substitution $\tau$ is a $W$-renaming if $\tau(x)$ is a variable for all $x\in W$ and $\tau$ is injective on $W$. Since $\sigma$ is the identity on $V\setminus\operatorname{Dom}(\sigma)$ and elements of $\operatorname{Dom}(\sigma)\setminus\operatorname{Dom}(\theta)$ are mapped to $\operatorname{Dom}(\sigma)\setminus\operatorname{Dom}(\theta)$ by $\sigma$, the substitution $\sigma$ is also a $W$-renaming for $W=V\setminus\operatorname{Dom}(\theta)$. By Lemma 2.12 in [28], there is a permutation $\pi$ that coincides with $\sigma$ on $W=V\setminus\operatorname{Dom}(\theta)$. Let $\sigma^{\prime}:=\pi^{-1}\sigma$. Then $\sigma^{\prime}\sim_{\emptyset}^{V}\sigma$, and thus also $\sigma^{\prime}\sim_{\mathcal{E\mkern-1.618muL}}^{V}\sigma$. In addition, for all $x\not\in\operatorname{Dom}(\theta)$ we have $\sigma^{\prime}(x)=\pi^{-1}\sigma(x)=x$, which yields $x\not\in\operatorname{Dom}(\sigma^{\prime})$. This shows that $\operatorname{Dom}(\sigma^{\prime})\subseteq\operatorname{Dom}(\theta)$. $\hfill\blacktriangleleft$\par
 
% END S3.SS0.SSS0.P0.SPx5.p4
 
% BEGIN S3.SS0.SSS0.P0.SPx5.p5
 
This lemma together with the next result shows that, up to equivalence, we can bound the variables occurring in more general substitutions by the substitutions they have as instance.\par
 
% END S3.SS0.SSS0.P0.SPx5.p5
 
% BEGIN Thmtheorem19
 
\par\medskip\noindent\textbf{Lemma 19.}\par
 
% BEGIN Thmtheorem19.p1
 
Let $\sigma,\theta$ be substitutions such that $\sigma\leq_{\mathcal{E\mkern-1.618muL}}^{V}\theta$ and $\operatorname{Dom}(\sigma)\subseteq\operatorname{Dom}(\theta)$. Then $\operatorname{VRan}(\sigma)\subseteq\operatorname{Dom}(\theta)\cup%
\operatorname{VRan}(\theta)$.\par
 
% END Thmtheorem19.p1
 
% END Thmtheorem19
 
% END S3.SS0.SSS0.P0.SPx5
 
% BEGIN S3.SS0.SSS0.P0.SPx6
 
\par\medskip\noindent\textbf{Proof.}\par
 
% BEGIN S3.SS0.SSS0.P0.SPx6.p1
 
Let $\lambda$ be such that $\lambda\sigma\approx_{\mathcal{E\mkern-1.618muL}}\theta$. Assume that $z\in\operatorname{VRan}(\sigma)$, but $z\notin\operatorname{Dom}(\theta)\cup\operatorname{VRan}(\theta)$. Then $z\not\in\operatorname{Dom}(\sigma)$, and thus $\lambda(z)=\lambda\sigma(z)\approx_{\mathcal{E\mkern-1.618muL}}\theta(z)=z$. However, since $z\in\operatorname{VRan}(\sigma)$, there is a variable $x\in\operatorname{Dom}(\sigma)$ such that $\sigma(x)$ contains $z$. Then $x\in\operatorname{Dom}(\theta)$, and thus all variables occurring in $\theta(x)$ belong to $\operatorname{VRan}(\theta)$. Since $\lambda\sigma(x)\approx_{\mathcal{E\mkern-1.618muL}}\theta(x)$, this means that $z\in\operatorname{VRan}(\theta)$, which contradicts our assumptions on $z$. $\hfill\blacktriangleleft$\par
 
% END S3.SS0.SSS0.P0.SPx6.p1
 
% BEGIN S3.SS0.SSS0.P0.SPx6.p2
 
The following lemma bounds the role depth (i.e., the maximal nesting of existential restrictions) as well as the concept constants and role names occurring in more general substitutions. Formally, the role depth $\operatorname{rd}(C)$ of an $\mathcal{E\mkern-1.618muL}$ concept description $C$ is defined inductively as follows:\par
  
\par\noindent  $\blacksquare$ 
% BEGIN S3.I1.i1.p1
 
$\operatorname{rd}(A):=\operatorname{rd}(\top):=0$ for all concept names $A$,\par
 
% END S3.I1.i1.p1
 
 
\par\noindent  $\blacksquare$ 
% BEGIN S3.I1.i2.p1
 
$\operatorname{rd}(C\sqcap D):=\max\{\operatorname{rd}(C),\operatorname{rd}(D)\}$ and $\operatorname{rd}(\exists r.C):=1+\operatorname{rd}(C)$.\par
 
% END S3.I1.i2.p1
 
  
% END S3.SS0.SSS0.P0.SPx6.p2
 
% BEGIN Thmtheorem20
 
\par\medskip\noindent\textbf{Lemma 20.}\par
 
% BEGIN Thmtheorem20.p1
 
Let $\sigma,\theta$ be substitutions such that $\sigma\leq_{\mathcal{E\mkern-1.618muL}}^{V}\theta$. Then the following holds for all $x\in V$: the role depth of $\sigma(x)$ is bounded by the role depth of $\theta(x)$, and the concept constants and role names occurring in $\sigma(x)$ also occur in $\theta(x)$.\par
 
% END Thmtheorem20.p1
 
% END Thmtheorem20
 
% END S3.SS0.SSS0.P0.SPx6
 
% BEGIN S3.SS0.SSS0.P0.SPx7
 
\par\medskip\noindent\textbf{Proof.}\par
 
% BEGIN S3.SS0.SSS0.P0.SPx7.p1
 
This is an easy consequence of the fact that $\approx_{\mathcal{E\mkern-1.618muL}}$ preserves the role depth as well as the set of concept constants and role names occurring in a concept description, and applying a substitution to a concept description can only increase the role depth and add concept constants or role names, but not decrease the role depth or remove concept or role names. $\hfill\blacktriangleleft$\par
 
% END S3.SS0.SSS0.P0.SPx7.p1
 
% BEGIN S3.SS0.SSS0.P0.SPx7.p2
 
As a consequence of the sequence of lemmas we have just shown we know that, for a given substitution $\theta$, the set of more general substitutions is finite up to the equivalence relation $\sim_{\mathcal{E\mkern-1.618muL}}^{V}$. In fact, Lemmas 18 and 19 show that one can restrict the attention to substitutions $\sigma$ satisfying $\operatorname{Dom}(\sigma)\cup\operatorname{VRan}(\sigma)\subseteq%
\operatorname{Dom}(\theta)\cup\operatorname{VRan}(\theta)$. In addition, for all $x\in\operatorname{Dom}(\sigma)$, the concept descriptions $\sigma(x)$ are built using role names and concept constants occurring in $\theta(x)$ as well as variables from $\operatorname{Dom}(\theta)\cup\operatorname{VRan}(\theta)$, and have a role depth that is bounded by the one of $\theta(x)$. It is well-known [13] that there are up to equivalence $\approx_{\mathcal{E\mkern-1.618muL}}$ only finitely many $\mathcal{E\mkern-1.618muL}$-concept descriptions satisfying these properties.\par
 
% END S3.SS0.SSS0.P0.SPx7.p2
 
% BEGIN Thmtheorem21
 
\par\medskip\noindent\textbf{Lemma 21.}\par
 
% BEGIN Thmtheorem21.p1
 
For a given substitution $\theta$, the set $\{\sigma\mid\sigma\leq_{\mathcal{E\mkern-1.618muL}}^{V}\theta\}$ is finite up to $\sim_{\mathcal{E\mkern-1.618muL}}^{V}$-equivalence.\par
 
% END Thmtheorem21.p1
 
% END Thmtheorem21
 
% BEGIN S3.SS0.SSS0.P0.SPx7.p3
 
As an immediate consequence we obtain that the set of substitutions more general than $\theta$ contains elements that are minimal w.r.t. $\leq_{\mathcal{E\mkern-1.618muL}}^{V}$. This obviously shows that the unrestricted unification type of $\mathcal{E\mkern-1.618muL}$ cannot be zero. Since the restricted unification type is zero, the unrestricted type cannot be unitary or finitary by Theorem 8.\par
 
% END S3.SS0.SSS0.P0.SPx7.p3

% BEGIN Thmtheorem22
 
\par\medskip\noindent\textbf{Theorem 22.}\par
 
% BEGIN Thmtheorem22.p1
 
The unrestricted unification type of $\mathcal{E\mkern-1.618muL}$ is infinitary.\par
 
% END Thmtheorem22.p1
 
% END Thmtheorem22

\begin{thebibliography}{99}

% BEGIN bib
 
\par\medskip\noindent\textbf{References}\par
  
\par\noindent  [1]  Majid Alizadeh, Mohammad Ardeshir, Philippe Balbiani, and Mojtaba Mojtahedi.  Unification types in Euclidean modal logics.  Log. J. IGPL, 31(3):422–440, 2023.  doi:10.1093/JIGPAL/JZAB036.  
 
\par\noindent  [2]  Franz Baader.  The theory of idempotent semigroups is of unification type zero.  J. Autom. Reason., 2(3):283–286, 1986.  doi:10.1007/BF02328451.  
 
\par\noindent  [3]  Franz Baader.  Characterization of unification type zero.  In Nachum Dershowitz, editor, Rewriting Techniques and Applications, 3rd International Conference, RTA-89, Proceedings, volume 355 of Lecture Notes in Computer Science, pages 2–14. Springer, 1989.  doi:10.1007/3-540-51081-8\_96.  
 
\par\noindent  [4]  Franz Baader.  Unification, weak unification, upper bound, lower bound, and generalization problems.  In Ronald V. Book, editor, Rewriting Techniques and Applications, 4th International Conference, RTA-91, Proceedings, volume 488 of Lecture Notes in Computer Science, pages 86–97. Springer, 1991.  doi:10.1007/3-540-53904-2\_88.  
 
\par\noindent  [5]  Franz Baader.  Unification in commutative theories, Hilbert’s basis theorem, and Gröbner bases.  J. ACM, 40(3):477–503, 1993.  doi:10.1145/174130.174133.  
 
\par\noindent  [6]  Franz Baader, Sebastian Brandt, and Carsten Lutz.  Pushing the $\mathcal{EL}$ envelope.  In Leslie Pack Kaelbling and Alessandro Saffiotti, editors, IJCAI-05, Proceedings of the Nineteenth International Joint Conference on Artificial Intelligence, pages 364–369. Professional Book Center, 2005.  URL: http://ijcai.org/Proceedings/05/Papers/0372.pdf.  
 
\par\noindent  [7]  Franz Baader and Wolfram Büttner.  Unification in commutative idempotent monoids.  Theor. Comput. Sci., 56:345–353, 1988.  doi:10.1016/0304-3975(88)90140-5.  
 
\par\noindent  [8]  Franz Baader and Silvio Ghilardi.  Unification in modal and description logics.  Log. J. IGPL, 19(6):705–730, 2011.  doi:10.1093/JIGPAL/JZQ008.  
 
\par\noindent  [9]  Franz Baader and Pierre Ludmann.  The exact unification type of commutative theories.  In Santiago Escobar and Mateu Villaret, editors, Proceedings of the 29th International Workshop on Unification (UNIF’15), pages 19–23, 2015.  URL: https://lat.inf.tu-dresden.de/research/papers/2015/BaLu-UNIF2015.pdf.  
 
\par\noindent  [10]  Franz Baader and Pierre Ludmann.  The unification type of ACUI w.r.t. the unrestricted instantiation preorder is not finitary.  In Silvio Ghilardi and Manfred Schmidt-Schauß, editors, Proceedings of the 30th International Workshop on Unification, UNIF 2016, pages 31–36, 2016.  URL: http://users.mat.unimi.it/users/ghilardi/UNIF2016/UNIF16-abstracts.pdf\#page=31.  
 
\par\noindent  [11]  Franz Baader, Carsten Lutz, and Boontawee Suntisrivaraporn.  CEL - A polynomial-time reasoner for life science ontologies.  In Ulrich Furbach and Natarajan Shankar, editors, Automated Reasoning, Third International Joint Conference, IJCAR 2006, Proceedings, volume 4130 of Lecture Notes in Computer Science, pages 287–291. Springer, 2006.  doi:10.1007/11814771\_25.  
 
\par\noindent  [12]  Franz Baader and Barbara Morawska.  Unification in the description logic $\mathcal{EL}$.  In Ralf Treinen, editor, Rewriting Techniques and Applications, 20th International Conference, RTA 2009, Proceedings, volume 5595 of Lecture Notes in Computer Science, pages 350–364. Springer, 2009.  doi:10.1007/978-3-642-02348-4\_25.  
 
\par\noindent  [13]  Franz Baader and Barbara Morawska.  Unification in the description logic $\mathcal{EL}$.  Log. Methods Comput. Sci., 6(3), 2010.  URL: http://arxiv.org/abs/1006.2289.  
 
\par\noindent  [14]  Franz Baader and Paliath Narendran.  Unification of concept terms in description logics.  J. Symb. Comput., 31(3):277–305, 2001.  doi:10.1006/JSCO.2000.0426.  
 
\par\noindent  [15]  Franz Baader, Thanh Binh Nguyen, Stefan Borgwardt, and Barbara Morawska.  Deciding unifiability and computing local unifiers in the description logic $\mathcal{EL}$ without top constructor.  Notre Dame J. Formal Log., 57(4):443–476, 2016.  doi:10.1215/00294527-3555507.  
 
\par\noindent  [16]  Franz Baader and Tobias Nipkow.  Term Rewriting and All That.  Cambridge University Press, 1998.  doi:10.1017/CBO9781139172752.  
 
\par\noindent  [17]  Franz Baader and Werner Nutt.  Combination problems for commutative/monoidal theories or how algebra can help in equational unification.  Appl. Algebra Eng. Commun. Comput., 7(4):309–337, 1996.  doi:10.1007/BF01195536.  
 
\par\noindent  [18]  Franz Baader and Jörg H. Siekmann.  Unification theory.  In Dov M. Gabbay, Christopher J. Hogger, J. A. Robinson, and Jörg H. Siekmann, editors, Handbook of Logic in Artificial Intelligence and Logic Programming, Volume 2, Deduction Methodologies, pages 41–126. Oxford University Press, 1994.  doi:10.1093/oso/9780198537465.003.0002.  
 
\par\noindent  [19]  Franz Baader and Wayne Snyder.  Unification theory.  In John Alan Robinson and Andrei Voronkov, editors, Handbook of Automated Reasoning, volume 1, pages 445–532. Elsevier and MIT Press, 2001.  doi:10.1016/B978-044450813-3/50010-2.  
 
\par\noindent  [20]  Leo Bachmair and Nachum Dershowitz.  Completion for rewriting modulo a congruence.  Theor. Comput. Sci., 67(2\&3):173–201, 1989.  doi:10.1016/0304-3975(89)90003-0.  
 
\par\noindent  [21]  Philippe Balbiani and Çiğdem Gencer.  About the unification types of modal logics.  In Nick Bezhanishvili, Rosalie Iemhoff, and Fan Yang, editors, Dick de Jongh on Intuitionistic and Provability Logics, pages 175–202. Springer International Publishing, 2024.  doi:10.1007/978-3-031-47921-2\_7.  
 
\par\noindent  [22]  Philippe Balbiani, Çigdem Gencer, Maryam Rostamigiv, and Tinko Tinchev.  Remarks about the unification types of some locally tabular normal modal logics.  Log. J. IGPL, 31(1):115–139, 2023.  doi:10.1093/JIGPAL/JZAB033.  
 
\par\noindent  [23]  Leonardo Manuel Cabrer and George Metcalfe.  From admissibility to a new hierarchy of unification types.  In Temur Kutsia and Christophe Ringeissen, editors, Proceedings of the 28th International Workshop on Unification, UNIF 2014, pages 41–46, 2014.  URL: https://www3.risc.jku.at/publications/download/risc\_5001/proceedings-UNIF2014.pdf\#page=46.  
 
\par\noindent  [24]  Leonardo Manuel Cabrer and George Metcalfe.  Exact unification and admissibility.  Log. Methods Comput. Sci., 11(3), 2015.  doi:10.2168/LMCS-11(3:23)2015.  
 
\par\noindent  [25]  Jacques Corbin and Michel Bidoit.  A rehabilitation of Robinson’s unification algorithm.  In R. E. A. Mason, editor, Proceedings of the 9th World Computer Congress, IFIP’83, pages 909–914, Paris, France, 1983. Elsevier.  
 
\par\noindent  [26]  Wojciech Dzik, Slawomir Kost, and Piotr Wojtylak.  Finitary unification in locally tabular modal logics characterized.  Ann. Pure Appl. Log., 173(4):103072, 2022.  doi:10.1016/J.APAL.2021.103072.  
 
\par\noindent  [27]  Wojciech Dzik, Slawomir Kost, and Piotr Wojtylak.  Unification types and union splittings in intermediate logics.  Ann. Pure Appl. Log., 176(1):103508, 2025.  doi:10.1016/J.APAL.2024.103508.  
 
\par\noindent  [28]  Elmar Eder.  Properties of substitutions and unifications.  J. Symb. Comput., 1(1):31–46, 1985.  doi:10.1016/S0747-7171(85)80027-4.  
 
\par\noindent  [29]  François Fages and Gérard P. Huet.  Complete sets of unifiers and matchers in equational theories.  In Giorgio Ausiello and Marco Protasi, editors, CAAP’83, Trees in Algebra and Programming, 8th Colloquium, Proceedings, volume 159 of Lecture Notes in Computer Science, pages 205–220. Springer, 1983.  doi:10.1007/3-540-12727-5\_12.  
 
\par\noindent  [30]  François Fages and Gérard P. Huet.  Complete sets of unifiers and matchers in equational theories.  Theor. Comput. Sci., 43:189–200, 1986.  doi:10.1016/0304-3975(86)90175-1.  
 
\par\noindent  [31]  Silvio Ghilardi.  Unification in intuitionistic logic.  J. Symb. Log., 64(2):859–880, 1999.  doi:10.2307/2586506.  
 
\par\noindent  [32]  Miki Hermann, Laurent Juban, and Phokion G. Kolaitis.  On the complexity of counting the Hilbert basis of a linear diophnatine system.  In Harald Ganzinger, David A. McAllester, and Andrei Voronkov, editors, Logic Programming and Automated Reasoning, 6th International Conference, LPAR’99, Proceedings, volume 1705 of Lecture Notes in Computer Science, pages 13–32. Springer, 1999.  doi:10.1007/3-540-48242-3\_2.  
 
\par\noindent  [33]  Michael Hoche, Jörg H. Siekmann, and Peter Szabó.  String unification is essentially infinitary.  FLAP, 3(5):755–788, 2016.  URL: http://www.collegepublications.co.uk/downloads/ifcolog00009.pdf.  
 
\par\noindent  [34]  Rosalie Iemhoff.  A syntactic approach to unification in transitive reflexive modal logics.  Notre Dame J. Formal Log., 57(2):233–247, 2016.  doi:10.1215/00294527-345997.  
 
\par\noindent  [35]  Emil Jerábek.  Blending margins: the modal logic K has nullary unification type.  J. Log. Comput., 25(5):1231–1240, 2015.  doi:10.1093/LOGCOM/EXT055.  
 
\par\noindent  [36]  Jean-Pierre Jouannaud and Claude Kirchner.  Solving equations in abstract algebras: A rule-based survey of unification.  In J.-L. Lassez and G. Plotkin, editors, Computational Logic: Essays in Honor of A. Robinson. MIT Press, Cambridge, MA, 1991.  
 
\par\noindent  [37]  Jean-Pierre Jouannaud and Helene Kirchner.  Completion of a set of rules modulo a set of equations.  SIAM J. Computing, 15(4):1155–1194, 1986.  doi:10.1137/0215084.  
 
\par\noindent  [38]  Donald E. Knuth and Peter B. Bendix.  Simple word problems in universal algebras.  In J. Leech, editor, Computational Problems in Abstract Algebra. Pergamon Press, Oxford, 1970.  doi:10.1016/B978-0-08-012975-4.50028-X.  
 
\par\noindent  [39]  Ralf Küsters.  Non-Standard Inferences in Description Logics, volume 2100 of Lecture Notes in Computer Science.  Springer, 2001.  doi:10.1007/3-540-44613-3.  
 
\par\noindent  [40]  Dallas Lankford.  Non-negative integer basis algorithms for linear equations with integer coefficients.  J. Autom. Reason., 5(1):25–35, 1989.  doi:10.1007/BF00245019.  
 
\par\noindent  [41]  Jean-Louis Lassez, Michael J. Maher, and Kim Marriott.  Unification revisited.  In Jack Minker, editor, Foundations of Deductive Databases and Logic Programming, pages 587–625. Morgan Kaufmann, 1988.  doi:10.1016/B978-0-934613-40-8.50019-1.  
 
\par\noindent  [42]  Alberto Martelli and Ugo Montanari.  An efficient unification algorithm.  ACM Trans. Program. Lang. Syst., 4(2):258–282, 1982.  doi:10.1145/357162.357169.  
 
\par\noindent  [43]  Mike Paterson and Mark N. Wegman.  Linear unification.  J. Comput. Syst. Sci., 16(2):158–167, 1978.  doi:10.1016/0022-0000(78)90043-0.  
 
\par\noindent  [44]  Gerald E. Peterson and Mark E. Stickel.  Complete sets of reductions for some equational theories.  J. of the ACM, 28(2):233–264, 1981.  doi:10.1145/322248.322251.  
 
\par\noindent  [45]  Gordon Plotkin.  Building in equational theories.  Machine Intelligence, 7:73–90, 1972.  
 
\par\noindent  [46]  Peter Raulefs, Jörg H. Siekmann, Peter Szabó, and E. Unvericht.  A short survey on the state of the art in matching and unification problems.  SIGSAM Bull., 13(2):14–20, 1979.  doi:10.1145/1089208.1089210.  
 
\par\noindent  [47]  John Alan Robinson.  A machine oriented logic based on the resolution principle.  J. of the ACM, 12(1):23–41, 1965.  doi:10.1145/321250.321253.  
 
\par\noindent  [48]  Manfred Schmidt-Schauß.  Unification under associativity and idempotence is of type nullary.  J. Autom. Reason., 2(3):277–281, 1986.  doi:10.1007/BF02328450.  
 
\par\noindent  [49]  Manfred Schmidt-Schauß.  A decision algorithm for distributive unification.  Theor. Comput. Sci., 208(1-2):111–148, 1998.  doi:10.1016/S0304-3975(98)00081-4.  
 
\par\noindent  [50]  Jörg H. Siekmann.  Matching under commutativity.  In Edward W. Ng, editor, Symbolic and Algebraic Computation, EUROSAM ’79, An International Symposium on Symbolic and Algebraic Computation, Proceedings, volume 72 of Lecture Notes in Computer Science, pages 531–545. Springer, 1979.  doi:10.1007/3-540-09519-5\_101.  
 
\par\noindent  [51]  Jörg H. Siekmann.  Universal unification.  In Robert E. Shostak, editor, 7th International Conference on Automated Deduction, CADE 7, Proceedings, volume 170 of Lecture Notes in Computer Science, pages 1–42. Springer, 1984.  doi:10.1007/978-0-387-34768-4\_1.  
 
\par\noindent  [52]  Jörg H. Siekmann.  Unification theory.  In Benedict du Boulay, David C. Hogg, and Luc Steels, editors, Advances in Artificial Intelligence II, Seventh European Conference on Artificial Intelligence, ECAI 1986, Proceedings, pages 365–400. North-Holland, 1986.  
 
\par\noindent  [53]  Jörg H. Siekmann.  Unification theory.  J. Symb. Comput., 7(3/4):207–274, 1989.  doi:10.1016/S0747-7171(89)80012-4.  
 
\par\noindent  [54]  Viorica Sofronie-Stokkermans.  Locality and applications to subsumption testing in $\mathcal{EL}$ and some of its extensions.  Sci. Ann. Comput. Sci., 23(2):251–284, 2013.  doi:10.7561/SACS.2013.2.251.  
 
\par\noindent  [55]  Viorica Sofronie-Stokkermans.  Representation theorems and locality for subsumption testing and interpolation in the description logics $\mathcal{EL}$, $\mathcal{EL}^{+}$ and their extensions with $n$-ary roles and numerical domains.  Fundam. Informaticae, 156(3-4):361–411, 2017.  doi:10.3233/FI-2017-1612.  
 
\par\noindent  [56]  Mark E. Stickel.  A unification algorithm for associative-commutative functions.  J. ACM, 28(3):423–434, 1981.  doi:10.1145/322261.322262.  
 
\par\noindent  [57]  Peter Szabó and Jörg H. Siekmann.  $E$-unification based on generalized embedding.  Math. Struct. Comput. Sci., 31(8):898–917, 2021.  doi:10.1017/S0960129522000019.  
 
\par\noindent  [58]  Katherine A. Yelick.  Combining unification algorithms for confined regular equational theories.  In Jean-Pierre Jouannaud, editor, Rewriting Techniques and Applications, First International Conference, RTA-85, Proceedings, volume 202 of Lecture Notes in Computer Science, pages 365–380. Springer, 1985.  doi:10.1007/3-540-15976-2\_18.  
  
% END bib
\end{thebibliography}
\end{document}
